\pagenumbering{roman}
\chapter*{Resumen}
\addcontentsline{toc}{chapter}{Resumen}
Este informe documenta una aplicación de tres en raya para Android, implementada en C\# con .NET MAUI. La persona juega con \texttt{X}; la computadora responde con \texttt{O} mediante una búsqueda Minimax que recorre las jugadas legales restantes, evalúa victorias, derrotas y empates y favorece ganar pronto o retrasar una derrota. El documento sigue la guía de laboratorio proporcionada para la asignatura: especifica requisitos funcionales y no funcionales, presenta diagramas de casos de uso, secuencia y clases, explica el diseño y describe el trabajo con una adaptación de SRUM+XP.

La revisión se apoya en el código del proyecto MAUI, en la configuración del proyecto, en tres capturas de Visual Studio Community 2026 y en las evidencias existentes del emulador. Estas últimas incluyen 14 imágenes y 25 verificaciones registradas. Los resultados documentan el turno de la IA, el bloqueo de una amenaza, una victoria de la computadora, un empate, el reinicio y la salida. También dejan visible que las opciones de los menús Archivo, Ver y Ayuda aparecen en pantalla pero no ejecutan las acciones anunciadas. El documento separa esos hallazgos de las mejoras propuestas y deja campos editables para los nombres y el grupo de trabajo.

\noindent\textbf{Palabras clave:} videojuegos, .NET MAUI, Android, C\#, Minimax, Tic-Tac-Toe, SRUM, XP.

\tableofcontents
\listoffigures
\listoftables
\clearpage
\pagenumbering{arabic}
